\documentclass[10pt,a4paper]{article}
\usepackage[margin=0.62in]{geometry}
\usepackage{fontspec}
\usepackage{xeCJK}
\usepackage[hidelinks]{hyperref}
\usepackage{enumitem}
\usepackage{tabularx}
\usepackage[normalem]{ulem}

\IfFontExistsTF{Microsoft JhengHei}{%
  \setCJKmainfont{Microsoft JhengHei}
}{%
  \IfFontExistsTF{PMingLiU}{%
    \setCJKmainfont{PMingLiU}
  }{%
    \setCJKmainfont{Noto Serif CJK TC}
  }%
}

\pagenumbering{gobble}
\setlength{\parindent}{0pt}
\setlength{\parskip}{0pt}
\setlist[itemize]{leftmargin=1.3em,itemsep=0pt,topsep=1pt,parsep=0pt,partopsep=0pt,label=\raisebox{0.2ex}{\tiny\textbullet}}
\setlist[enumerate]{leftmargin=1.5em,itemsep=1pt,topsep=0pt,parsep=0pt,partopsep=0pt}
\renewcommand{\arraystretch}{1.05}
\urlstyle{same}

\newcommand{\plainlink}[1]{\uline{\nolinkurl{#1}}}

\newcommand{\sectiontitle}[1]{%
  \vspace{5pt}
  {\large\textbf{#1}}\par
  \vspace{-1pt}
  \hrule
  \vspace{2pt}
}

\newcommand{\cventry}[3]{%
  \begin{tabularx}{\textwidth}{@{}Xr@{}}
    \textbf{#1} & \textit{#2} \\
    \multicolumn{2}{@{}X@{}}{#3}
  \end{tabularx}
  \par
}

\newcommand{\cventryshort}[2]{%
  \begin{tabularx}{\textwidth}{@{}Xr@{}}
    {\small\textbf{#1}} & {\small{#2}} \\
  \end{tabularx}
  \par
}

\newcommand{\projectentry}[2]{%
  \begin{tabularx}{\textwidth}{@{}>{\raggedright\arraybackslash}X@{}}
    {\small\textbf{#1}} \;|\; {\scriptsize #2}
  \end{tabularx}
  \par
}

\newcommand{\projectentrynolink}[1]{%
  \begin{tabularx}{\textwidth}{@{}>{\raggedright\arraybackslash}X@{}}
    {\small\textbf{#1}}
  \end{tabularx}
  \par
}

\newcommand{\educationentry}[2]{%
  \begin{tabularx}{\textwidth}{@{}Xr@{}}
    {\small\textbf{#1}} & {\small{#2}} \\
  \end{tabularx}
  \par
}

\begin{document}
\begin{center}
  {\LARGE\textbf{陳冠維}}\\[2pt]
  \small
  +886-975-518-289 \;|\;
  \href{mailto:elaping5691@gmail.com}{\uline{elaping5691@gmail.com}} \;|\;
  \href{https://github.com/balaboom123}{\uline{github.com/balaboom123}} \;|\;
  \href{https://www.linkedin.com/in/kwc-tw/}{\uline{linkedin.com/in/kwc-tw}} \;|\;
  \href{https://balaboom123.github.io/}{\uline{balaboom123.github.io}}
\end{center}

\sectiontitle{摘要}
國立中央大學資訊工程碩士生，專攻人工智慧，研究成果已發表於三場國際研討會。
具備以 PyTorch 與 TensorFlow 建構電腦視覺與自然語言處理系統之經驗，研究主題涵蓋場景文字辨識、手語翻譯與生物特徵辨識。
目標職缺為 AI/ML 工程師或電腦視覺研究相關職位。

\sectiontitle{學歷}
\educationentry
  {國立中央大學｜碩士，資訊工程學系，人工智慧組}
  {台灣  2025 -- 至今}
\vspace{1pt}
\educationentry
  {國立彰化師範大學｜學士，電機工程學系，(GPA 3.8/4.0)}
  {台灣  2021 -- 2025}

\sectiontitle{經歷}
\cventryshort
  {國立彰化師範大學 AI 應用實驗室｜\textit{專案助理}}
  {2024 年 6 月 -- 2024 年 12 月}
\begin{itemize}
  \item 與 8 位國際實習生於 TEEP 與 IIPP 研究計畫下合作，以 PyTorch 與 TensorFlow 完成 3 個生物特徵辨識與動作辨識之多模態 AI 原型，成果發表為 IET ICETA 同儕審查論文。
  \item 建立以 Pandas 與 NumPy 為基礎的前處理、資料集與評估管線，優化電腦視覺與語音相關研究任務中的模型實驗與驗證流程。
\end{itemize}

\sectiontitle{研究發表}
\begin{enumerate}[label={\arabic*.}]
  \item \textbf{K. Chen}, K. Zheng, and M. Tsai, ``STRLite: Lightweight Masked Autoencoders for Adaptable Scene Text Recognition,'' IEEE ICCE-TW（已接受），2026。
  \item \textbf{K. Chen} and T. Lin, ``SignDATA: Data Pipeline for Sign Language Translation,'' arXiv:2604.20357，2026。
  \item \textbf{K. Chen} and M. Tsai, ``Towards Compact Sign Language Translation: Frame Rate and Model Size Trade-offs,'' IEEE ICCE-TW（已接受），2026。
  \item \textbf{K. W. Chen}, T. Y. Lin, W. R. Yang, A. Kesarwani, and R. Singh, ``Two-step Authentication: Multi-biometric System Using Voice and Facial Recognition,'' IET Conf. Proc., 2024, doi: 10.1049/icp.2024.4141。
\end{enumerate}

\sectiontitle{專案}
\projectentry
  {STRLite：可調適場景文字辨識的輕量化遮罩自編碼器}
  {\plainlink{github.com/balaboom123/STRLite}}
\begin{itemize}
  \item 共同第一作者；設計 600 萬參數的 Vision Transformer（ViT）模型，結合 Masked Autoencoder（MAE）預訓練與自回歸解碼器微調，用於場景文字辨識。
  \item 於六個標準電腦視覺基準資料集上達成 93.82\% 加權準確率，並於 U14M 資料集上達成 81.03\%。
\end{itemize}
\vspace{3pt}
\projectentry
  {SignDATA：手語翻譯資料處理管線}
  {\plainlink{github.com/balaboom123/signdata-slt}}
\begin{itemize}
  \item 第一作者；開源專案累積 100+ stars；建構以設定驅動的資料前處理與資料集管理系統，用於標準化異質手語語料並支援 AI/ML 管線開發。
  \item 設計整合 YOLO/MMDet 偵測與 MediaPipe/MMPose 關鍵點擷取的處理流程，進行結構化特徵工程，並以系統設計原則支援多資料集擴充。
\end{itemize}
\vspace{3pt}
\projectentrynolink
  {邁向輕量化手語翻譯：幀率與模型規模之取捨}
\begin{itemize}
  \item 第一作者；建立 7,700 萬參數的端到端自然語言處理（NLP）管線。
  \item 在 BLEU 分數僅由 10.06 降至 9.53 的情況下，將運算量降低約 {\textasciitilde}75\%，提供模型效率取捨的實證分析。
\end{itemize}
\vspace{3pt}
\projectentry
  {雙步驟驗證：結合語音與人臉辨識的多生物特徵系統}
  {\plainlink{github.com/NCUE-EE-AIAL/Two-Step-Muti-Biometric-Authentication-System}}
\begin{itemize}
  \item 第一作者；帶領 2 位實習生設計結合機器視覺與音訊辨識／說話者驗證的多模態 AI/ML 系統。
  \item 設計基於 ResNet 的卷積神經網路（CNN），結合 Fbank 特徵擷取與 triplet loss 進行語音驗證（99.1\%）；並以 TensorFlow 微調結合 MTCNN 物件偵測的 VGG16 人臉辨識模型（95.1\%）。
\end{itemize}
\vspace{3pt}
\projectentry
  {SEO 內容智慧分析引擎}
  {\plainlink{github.com/balaboom123/SEO-Content-Intelligence-Engine}}
\begin{itemize}
  \item 以 TypeScript 建構全端 Next.js 網頁應用程式，爬取 Google SERP 前 10 名搜尋結果，透過 OpenAI 擷取 NLP 實體並進行分群，以逆向工程方式分析搜尋意圖，供 SEO 競爭分析使用。
  \item 實作 Supabase 身份驗證搭配 SQL 列級安全策略（Row-Level Security）、Recharts 互動式視覺化圖表、Google Sheets 匯出功能，並部署於 Vercel，加入速率限制與 API 安全強化。
\end{itemize}

\sectiontitle{技能}
\footnotesize
\begin{itemize}
  \item \textbf{程式語言與工具：} Python, C/C++, JavaScript, TypeScript, HTML/CSS, SQL, Git, Linux, Docker
  \item \textbf{ML 框架與技術棧：} PyTorch, TensorFlow, Scikit-learn, Hugging Face, OpenCV, NumPy, Pandas, YOLO, MMDet, MediaPipe
  \item \textbf{網頁開發：} React, Next.js, Vue, Supabase, Vercel
  \item \textbf{語言能力：} 英文（\href{https://github.com/balaboom123/balaboom123/blob/main/certifications/language/IELTS-2024-02-17.jpg}{\uline{IELTS 6.5}}、\href{https://github.com/balaboom123/balaboom123/blob/main/certifications/language/TOEIC-2023-07-27.jpg}{\uline{TOEIC 830}}）、中文（母語）、台語（母語）
\end{itemize}

\end{document}
